% Source: https://artofproblemsolving.com/community/c3h1668380_problem_about_modular_arithmetic
% Stable identity: aops-topic-1668380-q01
We need $x\equiv2\pmod5$, $x\equiv3\pmod7$ and $x\equiv6\pmod{12}$.

From the last condition, $x=12r+6$. Modulo $5$: $12r+6\equiv2r+1\equiv2$, so $2r\equiv1\equiv6$ and $r\equiv3\pmod5$. Writing $r=5s+3$ gives $x=60s+42$. Modulo $7$: $60s+42\equiv4s\equiv3$, and since $4\cdot6=24\equiv3$, $s\equiv6\pmod7$. Writing $s=7t+6$ gives
\[
\boxed{x=420t+402,\quad t\in\mathbb{Z}} .
\]
The smallest positive solution is $402$: indeed $402=5\cdot80+2=7\cdot57+3=12\cdot33+6$. Since $5$, $7$, $12$ are pairwise coprime, the Chinese remainder theorem confirms that the solutions form exactly one residue class modulo $5\cdot7\cdot12=420$.
